\chapter{Background Work and Busy UI}
\label{ch:background-busy}
\chapterquote{The user should see that work is happening; the EDT should not be the place where it happens.}{The busy rule}

\begin{chaptermap}
Core question & How do background operations interact with Swing screens safely?\\
Primary APIs & \api{SwingTaskServices}, \api{BusyOverlayBinding}, \api{BusyOverlayLayerUI}, \api{TaskService}, \api{TaskHandle}, \api{TaskCallbacks}, \api{CancellationToken}, and \api{GenerationCounter}.\\
Modern Java topic & Virtual-thread workers with explicit EDT completion.\\
Swing connection & Busy overlays, disabled commands, progress text, status binding, cancellation, and stale-result suppression.\\
Companion example & \modernJavaExample.\\
\end{chaptermap}

\section{Busy is state}

Busy state is not a paint trick. It is a presentation fact. It disables
commands, blocks incompatible input, shows overlays, changes cursor policy,
updates status text, controls cancellation, and prevents duplicate requests. If
busy state lives only as a local boolean inside a button listener, every other
part of the screen must guess.

\termdef{Busy state}{in Corusco} is normally a \api{ReadableValue<Boolean>}
exposed by a task service, task handle, presenter, or derived value. It can be
bound to buttons, commands, overlays, status labels, and tests.

\begin{lstlisting}[caption={Service-level busy state.}]
TaskService tasks = SwingTaskServices.virtualThreads();

ReadableValue<Boolean> busy = tasks.busy();
Binding overlay = BusyOverlayBinding.install(layer, busy);
\end{lstlisting}

The task service owns the worker execution. The overlay binding owns the visual
relationship. The screen scope owns both lifetimes.

\begin{figure}[htbp]
\centering
\begin{minipage}[t]{0.48\linewidth}
\includegraphics[width=\linewidth]{\busyTaskFigure}
\end{minipage}\hfill
\begin{minipage}[t]{0.48\linewidth}
\includegraphics[width=\linewidth]{\backgroundStatusFigure}
\end{minipage}
\caption{Busy work is seen in several places at once: overlays, disabled commands, cancellation, progress, and status text all observe modeled task state.}
\label{fig:background-busy-screens}
\end{figure}

This ownership sentence is more important than the particular classes. Swing
programs often fail here because "loading" is treated as a temporary detail of
one listener. The listener disables its button, starts a worker, and re-enables
the button later. Then a status line is added and receives its own flag. Then a
busy cursor is added around a different call. Then a cancel button appears and
must discover which operation is running. The screen has acquired a real piece
of state, but the code still treats it as scattered bookkeeping. Corusco pushes
that fact into a readable value, so the rest of the screen can observe it in a
uniform way.

The Swing module connects this model to concrete components; it does not make
background work magically safe. Core task contracts describe cancellation,
callbacks, handles, and generation counters without knowing about Swing. The
Swing task helpers add the worker-to-EDT delivery boundary, and busy overlay
bindings adapt the resulting state to visible input policy. This layered shape
is the same pattern used throughout Corusco: a stable model first, Swing
adaptation second, lifecycle ownership around both.

Busy state is therefore part of the screen's contract with the user. It says
that some operation is running and that certain other operations are either
temporarily unavailable, still available, or now have different meaning. A
Refresh button may be disabled while refresh is running. A Cancel button may be
enabled only while the matching task handle is busy. A detail form may remain
editable while the table refreshes, but a Save command may be disabled while a
save task is running. The same word "busy" hides several policies unless the
screen models them explicitly.

This is where many Swing applications become inconsistent. One listener sets a
wait cursor. Another disables a button. A third starts a \api{SwingWorker}. A
fourth writes a status label. A later feature adds an overlay. Each addition is
reasonable in isolation, but the screen no longer has one answer to the
question "what is running?" The user can sometimes click the same command
twice, sometimes edit hidden fields under an overlay, and sometimes watch an
old result overwrite new state. Corusco's task and busy APIs are small because
the hard part is not a large scheduler. The hard part is making the operation a
named, observable fact.

The first distinction is between operation state and presentation policy. A
\api{TaskService} can expose service-level busy state. A \api{TaskHandle} can
expose task-level busy state. A presenter can derive a screen-blocked value
from several task handles. A command can derive enabled state from dirty,
valid, and busy values. A busy overlay can observe one of those values and
block input. These are related facts, not synonyms. If the application uses one
global Boolean for all of them, it will eventually disable too much or too
little.

For example, consider a customer screen with a table on the left and an editor
on the right. Loading the table should probably block table interaction and
disable Refresh. It need not block editing the currently selected customer's
notes, if those notes are already loaded and save uses a separate task.
Saving the notes should disable Save and perhaps the notes editor, but it need
not block scrolling the table. Loading details for a newly selected customer
should suppress stale detail results, but it need not cancel a background table
refresh. A single \api{busy} flag cannot express this. Several named busy facts
can.

The second distinction is between "work is running" and "input is incompatible."
A screen may show that validation is running while still allowing the user to
continue typing. A search field may remain editable while a previous query is
running, with generation counters suppressing stale results. A destructive
Delete command may be disabled while any save task is running. The overlay is
only one presentation of one policy: this region should not accept input while
this busy value is true. Do not use an overlay as a substitute for deciding
which input is incompatible.

The third distinction is between visible busy and lifecycle busy. A task
started by a dialog may still be running when the dialog closes. The user may
not need a cancellation message, because the screen is gone. But the task
still needs cancellation or stale-result suppression; otherwise it may publish
into a model whose view no longer exists. A task started by an application-wide
sync service may outlive a single window and should not be cancelled merely
because one view closed. The owner of the task decides the lifecycle. The view
decides which part of that lifecycle is visible.

The fourth distinction is between responsiveness and speed. Moving work off
the EDT does not make the work faster. It keeps the UI event queue available
for painting, input, cancellation, and status updates while work proceeds. A
slow database query may still take two seconds. The user's experience changes
because the screen can show progress, accept cancellation, and avoid "not
responding" behavior. If the callback then performs another two-second query
on the EDT, the benefit is lost. The boundary must stay clean all the way to
result application.

\begin{principlebox}{Busy Is Not a Spinner}
Treat busy as modeled interaction policy: which operation is running, who owns
it, what can be cancelled, what input is incompatible, what result is current,
and how success, failure, and cancellation are shown.
\end{principlebox}

A useful design exercise is to write the busy policy before writing the task
code. What starts the operation? Can it be started twice? What becomes disabled?
What remains usable? Does the operation block a component subtree, a command
family, one table row, or the whole screen? Is cancellation visible? What
happens to old data while new data loads? What happens if a newer request
finishes before an older request? What happens if the owner closes? These
questions turn "run it in the background" into a real UI contract.

The answer should appear in model state. A task handle belongs to one submitted
operation. A task service belongs to a presenter, dialog, or application
lifecycle. A generation counter belongs to a request stream whose old results
must be ignored. A progress model belongs to an operation that needs more than
a Boolean. Commands derive enabled state from these facts. Overlays observe
these facts. Tests assert these facts. Once the model exists, Swing becomes
the adapter again rather than the storage location for concurrency truth.

\textbf{Task execution and EDT delivery.}

Background work may use platform threads, virtual threads, executors, task
services, or legacy \api{SwingWorker}. The Swing rule is unchanged: component
mutation belongs on the EDT. A virtual thread can load data; it must not call
\api{JTable.setModel} directly. A callback can update Swing only when the task
service promises EDT delivery.

\api{SwingTaskServices.virtualThreads} creates a virtual-thread-backed
\api{TaskService} whose terminal callbacks and final busy-state updates are
delivered through \api{SwingEdt.executor}. Task bodies run away from the UI.
Callbacks return to Swing.

Virtual threads make this style more attractive on modern Java, but they do not
relax Swing's rule. A virtual thread is cheap to block while a repository waits
for I/O; it is not a license to mutate a component from outside the EDT. Treat
the task body as a computation that produces data and the callback as the
short, EDT-confined handoff that applies data to presentation models. The more
clearly this line is drawn, the less the program depends on timing accidents.

Structured concurrency, where available in the surrounding application, can
help organize related non-UI subtasks. The boundary to Swing remains explicit.
Gather data concurrently away from the EDT, combine it into a result value, and
then deliver that result through a Swing-aware callback. Do not pass live Swing
components into a structured task scope. The object crossing the boundary should
be data, not a widget.

\begin{lstlisting}[caption={Virtual-thread work with Swing callback delivery.}]
TaskService tasks = SwingTaskServices.virtualThreads();

tasks.submit(
        cancellation -> {
            cancellation.throwIfCancellationRequested();
            return repository.loadCustomers(cancellation);
        },
        TaskCallbacks.onSuccess(rows -> {
            tableRows.batch(list -> {
                list.clear();
                rows.forEach(list::add);
            });
        }));
\end{lstlisting}

The success callback is short. It applies the already-produced result. It does
not perform another slow query on the EDT.

\textbf{What belongs in the task body.}

Put unpredictable work in the task body: I/O, database calls, service calls,
file reads, parsing, report generation, and other operations whose latency is
not bounded by a tiny constant. The task body receives a
\api{CancellationToken}. Check it at useful boundaries.

\begin{lstlisting}[caption={Task body with cancellation boundaries.}]
UiTask<CustomerReport> task = cancellation -> {
    cancellation.throwIfCancellationRequested();
    Customer customer = repository.loadCustomer(id);
    cancellation.throwIfCancellationRequested();
    List<Order> orders = repository.loadOrders(id);
    cancellation.throwIfCancellationRequested();
    return reportFactory.create(customer, orders);
};
\end{lstlisting}

Do not update components inside this body. Return data.

This advice includes indirect component updates. A task body should not mutate
an observable list that is already bound to a \api{JTable} unless that list has
an explicit EDT adaptation that makes the operation safe. It should not write
to a form field model if the form is immediately bound to Swing and the binding
expects synchronous EDT notifications. A background task can build a new list,
read a file, parse a document, or call a service. The callback can then replace
the observable state on the EDT in one clear step.

\textbf{What belongs in callbacks.}

Callbacks should be short and deterministic. They update model values, row
sources, status text, or command state. They may show a small error message if
the application policy says so. They should not block.

\api{TaskCallbacks} has terminal paths for success, failure, and cancellation.
Only one path should run. Cooperative cancellation should suppress success.
Failure should be visible enough for diagnostics and humane enough for users.

A good callback is easy to read because it has no hidden second job. It does
not start a new slow query, wait for a lock, or perform an export. If the
success path receives a thousand rows, it should apply them with a batched list
mutation. If it receives a domain object, it should set a value or refresh a
form. If it receives a failure, it should translate that failure into the
screen's problem, status, or error policy. The callback is not the task; it is
the delivery receipt.

\textbf{Task handles.}

\api{TaskHandle<T>} observes one submitted task. It exposes task-level busy
state, cancellation, and completion. Closing the handle requests cancellation.
Register handles with the view or dialog lifecycle when the task should stop
when the screen closes.

\begin{lstlisting}[caption={A view scope owns task handles.}]
TaskHandle<CustomerReport> handle = tasks.submit(task, callbacks);
viewScope.add(handle);

cancelButton.addActionListener(event -> handle.cancel());
\end{lstlisting}

The service may outlive one task. The handle belongs to the operation.

\textbf{Service busy versus task busy.}

Service-level busy means at least one task owned by the service is running.
Task-level busy means one particular operation is running. Use the right one.

A whole-screen overlay often observes service busy. A progress line beside a
single row should observe a task handle. A Refresh command may be disabled while
the refresh handle is busy but remain independent of a background validation
task. When commands become too entangled, derive specific busy values rather
than using one global flag for everything.

\textbf{Local busy scopes.}

Some screens need several busy regions. A left table may refresh while a right
detail form remains editable. A background validation task may mark one field
busy without blocking the entire dialog. A report export may disable only export
commands. Model these scopes explicitly.

\begin{lstlisting}[caption={Separate busy facts for separate policies.}]
ReadableValue<Boolean> tableBusy = refreshTask.busy();
ReadableValue<Boolean> saveBusy = saveTask.busy();

DerivedValue<Boolean> formBlocked = DerivedValue.of(
        () -> Boolean.TRUE.equals(saveBusy.value()),
        saveBusy);
\end{lstlisting}

Avoid the single global busy flag unless the whole screen truly has one
blocking state.

This is a common source of awkward user experience. A screen that disables
everything for every operation feels simpler to implement, but it teaches users
that the program is less capable than it is. Conversely, a screen that leaves
everything enabled during every operation invites duplicate requests and stale
state. The correct answer is usually neither global blocking nor complete
freedom. Name the region or command family that is affected, model that busy
fact, and bind only the incompatible interactions to it.

\section{Task lifecycle and worker-to-EDT handoff}

A \api{TaskService} has two responsibilities that should not be confused. It
runs task bodies on a worker executor, and it delivers terminal callbacks
through a configured callback executor. The core task API does not say "EDT";
it says "callback executor." The Swing factory chooses \api{SwingEdt.executor}
for that callback executor. This is the small seam where toolkit-neutral tasks
become Swing-safe task delivery. It is also the line a reviewer should trace
when a background result mutates the UI.

The default virtual-thread service owns its worker executor. Closing the
service cancels outstanding handles and shuts down the worker executor it owns.
A service created over caller-owned executors cancels outstanding handles but
does not close the caller's executor service. This is not a minor implementation
detail. It decides who is responsible for resource cleanup. A dialog-local
virtual-thread service can be closed with the dialog. An application-wide
executor belongs to the application and should not be shut down by a single
panel.

Service-level busy is implemented as an aggregate count of active tasks. The
first active task sets busy true. The last completed, failed, or cancelled task
sets busy false through the callback executor. In a Swing service, the final
busy false event therefore arrives on the EDT. A screen-level overlay that
observes service busy will remain active while any service task is active. If
two unrelated operations use the same service and one should not block the
other's screen region, that is a design smell: use separate services, separate
task handles, or derived busy values.

A \api{TaskHandle} is the operation-level object. It has its own busy value,
cancellation token, completion future, cancel method, and \api{close} method
that delegates to cancellation. This is the object a Cancel button usually
needs. It is also the object a view scope should own when the operation should
stop with the view. The service may outlive the view; the handle usually does
not. Keeping this distinction sharp prevents both leaks and over-cancellation.

Cancellation is cooperative even when the underlying future is cancelled. The
service signals the token and attempts to cancel the worker future, but task
code must still check the token at useful boundaries. A blocking library call
may not notice cancellation immediately. A CPU loop will not stop unless it
checks. A parser may need to check between records. A repository method that
accepts a cancellation token should pass it into lower layers or at least check
between expensive calls. The UI can request cancellation; task code must make
that request effective.

\begin{lstlisting}[caption={A task service registered with a view lifecycle.}]
final class CustomerPresenter implements AutoCloseable {
    private final TaskService tasks = SwingTaskServices.virtualThreads();
    private TaskHandle<List<CustomerRow>> refresh;

    void refresh() {
        if (refresh != null && !refresh.isDone()) {
            return;                         // Policy: one refresh at a time.
        }
        refresh = tasks.submit(
                cancellation -> repository.loadRows(cancellation),
                TaskCallbacks.onSuccess(rows -> tableRows.replaceAll(rows)));
    }

    void cancelRefresh() {
        if (refresh != null) {
            refresh.cancel();
        }
    }

    @Override
    public void close() {
        tasks.close();
    }
}
\end{lstlisting}

The listing shows a conservative policy: one refresh at a time. Another screen
might allow overlapping searches and use a generation counter to accept only
the newest result. The task API supports both policies because the policy is
not hidden inside the executor. The presenter decides what should happen when
the user asks for another operation before the first one finishes.

Terminal callbacks are exclusive. A successful task calls success. An
unexpected exception calls failure. A cooperative cancellation calls
cancellation and suppresses success/failure. This exclusivity is what lets UI
code reason about final state. A task should not both succeed and then be
treated as cancelled because a button was pressed late. A cancelled task should
not publish a success result merely because the repository returned after the
token was requested. The service checks cancellation before and after the task
body, and task code should check inside the body at longer boundaries.

Completion futures are useful in tests and in non-UI orchestration, but they
should not tempt UI code into blocking. Calling \api{completion().get()} on the
EDT defeats the entire task model. Tests can wait with timeouts because they
control the scenario. Production Swing callbacks should be expressed as
callbacks, value updates, command state, and model changes. If another task
must begin after one completes, start it from the callback or model transition,
not by blocking the event queue.

Failure handling should preserve diagnostic evidence. A failure callback can
translate the exception into a problem set, status message, error dialog, or
support log. It should not merely print a stack trace from a worker thread.
Silent background failure is worse than a foreground exception because the
user sees stale or missing data with no explanation. At minimum, the screen
should move from busy to a visible failed state and the exception should remain
available for diagnostics.

Virtual threads are a good fit for UI-initiated I/O because they make blocking
task bodies cheap. They do not make Swing components thread-safe, do not remove
the need for cancellation checks, and do not decide result ordering. A virtual
thread can call a blocking repository method. It should return data. The EDT
callback should apply data. If a virtual thread holds a lock needed by the EDT
or calls back into Swing indirectly, the program can still freeze. Modern Java
reduces thread-management friction; it does not cancel the Swing contract.

Structured concurrency has a similar boundary. A task body may use structured
subtasks to load customer, orders, and preferences concurrently, then combine
them into one immutable result. That structured scope should live away from
Swing. The result crossing into the callback should be data. If the structured
scope receives a \api{JTable}, \api{JTextField}, or mutable Swing model, the
boundary has been breached. Keep Swing mutation in the callback and keep the
callback short.

Close paths deserve tests. Closing a task service should reject later
submissions. Closing a task handle should request cancellation. Closing a view
scope should close overlay bindings and task handles. Closing while a callback
is queued should not reinstall visual state. These are race-prone paths that
manual testing rarely exercises because the timing window is small. Automated
tests can make the timing window explicit with latches, controlled executors,
and EDT drains.

\section{Modern Java boundaries for Swing tasks}

Modern Java gives Swing programmers better background tools, but it does not
change Swing's event-thread rule. Virtual threads make it cheap to write a task
body that blocks while waiting for a database, file system, or remote service.
Records make task results easier to name and pass across the boundary. Pattern
matching makes result handling and error classification clearer. Structured
concurrency can organize related subtasks inside a task body. None of these
features should move component mutation away from the EDT.

The safest modern-Java shape is data-in, data-out. A command captures the
current model snapshot on the EDT. The task body receives immutable input and
returns an immutable result, often a record. The callback applies the result on
the EDT. If the task fails, the callback translates the failure into modeled
screen state. If cancellation is requested, the task body stops cooperatively
and the cancellation callback applies the screen's cancellation policy. This
shape works with virtual threads, platform executors, or test executors because
the Swing boundary is explicit.

\begin{lstlisting}[caption={Records make task input and output explicit.}]
record ExportRequest(Path file, List<CustomerRow> rows) {
}

record ExportResult(Path file, int rowCount, Duration elapsed) {
}

ExportRequest request = new ExportRequest(targetFile, List.copyOf(tableRows));
TaskHandle<ExportResult> handle = tasks.submit(
        cancellation -> exporter.export(request, cancellation),
        TaskCallbacks.onSuccess(result ->
                status.setValue("Exported " + result.rowCount() + " rows")));
\end{lstlisting}

The request is a snapshot. It does not contain a \api{JTable}. It does not
contain a live row sorter. It does not read from Swing while the worker runs.
The result is also data. The callback can safely translate it into status text,
problem state, command state, or a support event. Records are not required, but
they make the boundary visible in code review.

Virtual threads should be used for the kind of work they are good at: many
blocking operations whose cost is mostly waiting. They are not magic for CPU
bound rendering, geometry simplification, image processing, compression, or
encryption. CPU-heavy work still competes for cores. A virtual thread running a
large CPU loop can starve other work just as a platform thread can. For CPU
heavy operations, use bounded executors, chunk work, check cancellation, and
avoid flooding the EDT with progress. The task abstraction remains the same,
but executor choice is a performance policy.

Structured concurrency is useful inside a task body when several subtasks
belong to one result. A detail load may need customer, orders, and permissions.
A report export may need data, template, and resource bundle. A map layer may
need several tiles. A structured scope can start those subtasks together and
cancel siblings if one fails. The callback should still receive one combined
result or failure. Do not let several subtasks race independently to update
Swing components. That would recreate stale-result bugs under a newer syntax.

Pattern matching can make failure policy clearer when the application defines
failure types. A repository exception, validation exception, cancellation, and
authorization failure should not all become the same dialog. A failure callback
can classify the exception and update the right model. Keep this classification
short on the EDT. If classification requires slow I/O, it belongs in the task
body or in a follow-up task, not in the callback.

Modern Java also makes it easier to accidentally capture too much. A lambda
task body can close over \api{this}, a presenter, a component, or a mutable
list without any ceremony. Review task bodies for captured Swing state. If a
lambda references \api{nameField}, \api{table}, \api{panel}, or a mutable
observable model that is bound to Swing, ask whether it should instead capture
a snapshot. Concise syntax is useful only when the captured boundary is still
honest.

Modules and packaged runtimes add another operational detail. Background tasks
that load resources, templates, fonts, or service providers must work inside
the packaged image, not only in the IDE. A task that exports a report should
load resources through module-aware resource APIs. A task that writes a file
should handle permissions and user-selected paths. Busy UI is not a substitute
for packaging correctness; it merely gives the user a coherent experience while
the packaged operation runs.

The rule for this chapter is therefore conservative and modern at the same
time. Use virtual threads where they simplify blocking work. Use records to
make inputs and outputs explicit. Use structured concurrency inside task
bodies when subtasks belong to one operation. Use pattern matching to classify
results and failures clearly. Keep Swing mutation on the EDT. Keep callbacks
short. Keep stale-result and cancellation policy explicit. Modern Java should
make the old Swing rules easier to obey, not easier to forget.

\section{Busy presentation}

\api{BusyOverlayBinding} connects an observable busy value to a \api{JLayer}
overlay. It installs a \api{BusyOverlayLayerUI}, forwards initial and later busy
values, schedules visual updates onto the EDT when needed, and restores the
previous layer UI when closed.

\begin{lstlisting}[caption={Wrapping a form in a busy layer.}]
JLayer<JComponent> layer = new JLayer<>(formPanel);

Binding busyOverlay = BusyOverlayBinding.install(layer, presenter.busy());
scope.add(busyOverlay);
\end{lstlisting}

The overlay is more than a translucent rectangle. While busy,
\api{BusyOverlayLayerUI} consumes mouse, mouse-motion, mouse-wheel, key, and
focus input events before they reach the wrapped view. It also switches the
cursor to a wait cursor and repaints. This is input policy.

Readers who have written only web applications may think of an overlay as
purely visual. In Swing, a \api{JLayer} overlay can be a disciplined input
boundary. It sits around the component tree and can decide whether events pass
through to the wrapped view. That is why a busy overlay belongs in the same
conversation as command enablement. Both are policies about what the user may
do while work is in progress. One policy acts through actions and buttons; the
other acts through the event stream reaching the component subtree.

\begin{pitfallbox}{Overlay Without Policy}
A decorative overlay that does not block incompatible input may make the screen
look busy while still allowing duplicate edits. Decide whether the operation
should block input, and implement that policy explicitly.
\end{pitfallbox}

The \api{JLayer} choice is important. It lets the overlay sit around an
existing component subtree without replacing the subtree's components. The
wrapped form still has its fields, layout, validation, and focus state. The
layer UI paints after the wrapped view and can intercept events before they
reach it. This is very different from disabling every child component. Disabled
children may change appearance, lose focus behavior, and require a traversal of
the component tree. A layer can express "the region is temporarily blocked"
without rewriting every child state.

That does not mean a layer is always the right policy. If a save operation
should disable only the Save command while the rest of the form remains
editable, an overlay is too broad. If a destructive import should prevent all
interaction with a table while leaving a Cancel command active outside the
layer, wrapping only the table is appropriate. If a global database migration
blocks the entire application window, a root-level glass pane or top-level
layer may be appropriate. Choose the overlay boundary from the interaction
policy, not from where it is easiest to put a panel.

\api{BusyOverlayBinding} installs the layer UI immediately and restores the
previous UI on close. This matters in real applications because a \api{JLayer}
may already have a UI for validation marks, debugging, or another visual
effect. A binding that simply overwrote the UI and never restored it would
destroy the earlier behavior. The busy binding snapshots the previous UI,
installs the busy UI, forwards current and future busy values, and restores the
previous UI on close. It owns the temporary relationship, not the layer itself.

Busy value delivery may come from outside the EDT. The binding therefore
schedules visual updates onto the EDT and ignores queued updates after close.
This is a small but essential contract. Without it, a background value change
could call \api{setBusy} on a Swing UI delegate from a worker thread, or a late
queued update could reactivate an overlay after the view closed. The binding is
not just a convenience call; it is the dispatcher and lifecycle adapter between
observable busy state and Swing component state.

The layer UI itself is deterministic. It paints a translucent cover with
\api{AlphaComposite.SrcOver}, uses a wait cursor while busy, and consumes mouse,
mouse-motion, mouse-wheel, key, and focus input events while busy. It avoids
animation. This is a deliberate documentation and testing choice. A deterministic
overlay can be captured in screenshots, tested with a small offscreen
image, and reasoned about from state alone. Applications can add animated
progress elsewhere when product design requires it.

Input consumption by an overlay should be coordinated with command enablement.
If the overlay consumes input over the form but a toolbar Save command remains
enabled outside the layer, the user may still trigger an incompatible operation.
If command enablement disables Save but a field inside the layer still accepts
typing because there is no overlay or component disablement, state may change
under a running operation. Visual blocking and command blocking should observe
the same modeled policy even though they act through different Swing
mechanisms.

Cursor policy is also state. A wait cursor says "this region is busy." If the
operation does not block that region, a wait cursor over it may be misleading.
If a task is running in the background but the user can continue editing, a
status message or small progress indicator may be better than a wait cursor.
If a region is blocked, the wait cursor should be restored when the binding
closes, even if the busy value changes later. This is the same restoration
discipline used for borders and tooltips.

Accessibility should not be forgotten. A visual overlay alone may not explain
why keyboard input is blocked or what operation is running. A busy screen can
update status text, accessible descriptions, or command disabled reasons from
the same busy model. If cancellation is available, the Cancel command should be
reachable by keyboard and described as such. The overlay may consume input in
one subtree, but it should not make the whole application inaccessible without
another path to cancellation or status.

\textbf{Deterministic overlay rendering.}

The standard overlay is intentionally animation-free. Animations can be useful,
but deterministic overlays are easier to test and screenshot. A book example or
generated screenshot should not fail because an animation frame landed at a
different moment.

If an application wants animated progress, add it deliberately and keep tests
aware of it. The default busy overlay should communicate blocking state without
becoming a miniature rendering engine.

\textbf{Disabled commands.}

Busy state often disables commands. But "busy" is rarely the only reason a
command is disabled. Save may be disabled because the form is clean, because
there are validation errors, or because a save is in progress. Refresh may be
disabled because a refresh is in progress but still allowed while async
validation is running.

Use derived values to state the policy:

\begin{lstlisting}[caption={Command enablement derived from busy and form state.}]
DerivedValue<Boolean> canSave = DerivedValue.of(
        () -> dirty.value()
                && !problems.value().hasErrors()
                && !saveBusy.value(),
        dirty, problems, saveBusy);

canSave.subscribe(event ->
        saveCommand.setEnabled(Boolean.TRUE.equals(event.newValue()),
                event.origin()));
\end{lstlisting}

The command observes policy. The button observes the command.

\textbf{Progress and status.}

The current task service exposes busy state, not a full progress model. For
operations that need progress, model progress explicitly: current step, maximum,
message, cancellable flag, and perhaps estimated remaining time. Do not squeeze
all of this into a status label string.

Status text should be derived from task state. "Loading customers..." belongs to
the load task. "Loaded 240 customers" belongs to success. "Load cancelled"
belongs to cancellation. "Load failed" belongs to failure. This makes tests and
translation easier.

\textbf{Long-running streams.}

Some work is not a single request with one result. A log tail, event stream,
large import, or live calculation may deliver many intermediate updates. The
same principles apply, but the model should expose an observable list, progress
value, or event stream rather than pretending there is one final result.

Batch intermediate row updates. Avoid publishing thousands of single-row events
when the UI only needs periodic refresh. Use cancellation to stop the stream and
lifecycle scopes to close subscriptions. If the stream can outlive the view,
the view must detach cleanly.

\textbf{Back-pressure and coalescing.}

Desktop applications often receive background updates faster than humans can
see them. A table that receives one event for every imported row may spend more
time repainting than helping the user. Coalesce updates where possible. Batch
list mutations. Debounce status messages. Persist table state with a scheduler
rather than writing preferences during every resize pixel.

Coalescing is not lying about state. It is choosing a truthful update boundary
that the UI can handle.

The same principle applies to progress reporting. A progress bar that moves ten
thousand times per second is not more accurate for a human user than one that
moves at visible milestones. A status label that displays every imported file
name may be less useful than one that displays the current batch and a count.
Choose update boundaries that preserve semantic truth while giving Swing enough
time to lay out, repaint, and process input.

Progress deserves its own model when it affects more than one string. A useful
progress model may include a phase name, current count, maximum count,
indeterminate flag, message, cancellable flag, and severity. A simple task can
publish only busy true/false. A file import probably needs more. A report
export may need "preparing", "rendering", "writing", and "done" phases. A
background validation sweep may need a field count and a current problem
summary. Once progress has structure, a label, progress bar, tooltip, status
line, and support log can all derive from the same state.

Do not confuse progress precision with progress usefulness. A repository may
know exact record counts. A remote service may not. A task may have phases
where the maximum is known and phases where it is not. A progress bar should
not lie by showing precise percentages for indeterminate work. Use an
indeterminate state when the maximum is unknown; switch to determinate when it
becomes known. The state transition is part of the task model, not a random
choice made by the progress bar.

Streams and imports need batching rules before they need fancy rendering. If an
import reads 100,000 rows and publishes each row to an observable list on the
EDT, the table may become the bottleneck even though file reading is off the
EDT. Better designs accumulate a batch away from the EDT and publish a bounded
batch on the EDT, or publish into a staging model and let the table refresh at
human cadence. This is not merely performance tuning; it protects input
latency and repaint stability.

The batch size is a policy. A small batch makes progress feel immediate but
causes more UI events. A large batch reduces event overhead but may make the UI
feel frozen between updates. A time-based policy can work well: publish at most
every 50 or 100 milliseconds while keeping a final flush at completion. Tests
can verify that the importer eventually publishes all rows and that cancellation
stops publication. Performance tests can tune the cadence for real data.

Coalescing should happen near the boundary where high-frequency events become
UI state. A parser can report exact internal progress to a task-local counter.
The presenter can coalesce that into a readable progress value. The binding can
update the progress bar. If coalescing is hidden inside the progress bar, other
surfaces such as status labels and support diagnostics may still receive too
many updates. Model first; Swing adapter second.

Stale-result suppression and cancellation solve different problems. Cancellation
asks old work to stop. Generation counters prevent old work from publishing if
it finishes anyway. A well-behaved repository may stop promptly when cancelled,
but a remote call may complete despite cancellation. A generation check in the
callback is still useful. Conversely, a generation check does not reduce load
on the remote service or local CPU; cancellation is still useful. Mature
screens often use both.

The generation should be advanced when the user action makes older results
obsolete, not when the old task happens to finish. In a search field, advance
when a new search request is scheduled. In a master-detail screen, advance
when selection changes and a new detail load starts. In a reset path, call
\api{invalidate} so queued callbacks cannot repopulate cleared state. The
token is captured with the task and checked in the callback before applying
the result. The result consumer still has to run on the appropriate executor.

\begin{lstlisting}[caption={Search result suppression with cancellation and generation.}]
void search(String text) {
    if (currentSearch != null) {
        currentSearch.cancel();
    }
    GenerationCounter.Generation generation = generations.advance();
    currentSearch = tasks.submit(
            cancellation -> repository.search(text, cancellation),
            TaskCallbacks.onSuccess(rows ->
                    generations.tryAccept(generation, rows,
                            resultRows::replaceAll)));
}
\end{lstlisting}

This example intentionally cancels the old task and still checks the generation
in the success callback. The cancellation request is a resource policy. The
generation check is a state policy. If the old task ignores cancellation and
finishes, its result is stale. If the old task stops promptly, the generation
check is cheap. The screen is correct in both cases.

Failure policy should be proportional to the operation. A failed background
refresh may leave old data visible with a status message. A failed save may
need a dialog, a problem set, and preserved dirty state. A failed optional
preview may only mark that preview unavailable. A failed export may need a
file-path-specific message and a support log entry. Do not build one generic
"background task failed" dialog and use it everywhere. The task service reports
failure; the presenter decides what failure means in the workflow.

Cancellation policy is equally contextual. Cancelling a search usually means
"ignore that search and let me continue typing." Cancelling an import may mean
"keep imported rows so far" or "rollback the import"; the product must decide.
Cancelling a save may not be possible once the server has committed the change.
Cancelling because a dialog closed may need no visible message. The
\api{cancelled} callback is a terminal path, but the user-facing state after
that path is a workflow decision.

Close paths combine cancellation and generation. When a view closes, cancel
task handles owned by the view and invalidate generation counters whose results
would update that view. Close overlay bindings so late busy values do not
reinstall the layer UI. Close subscriptions so progress values no longer update
dead labels. A close path that only hides the window but leaves tasks and
bindings alive is a delayed bug; it may appear as a leak, a stale update, or an
exception in a component that is no longer displayed.

\textbf{Stale results, failure, and cancellation.}

The most dangerous task bug is not failure. It is an old success result that
arrives after the user selected a different record and silently overwrites newer
state. Cancellation helps, but it is not sufficient. Use generation counters for
request streams.

\begin{lstlisting}[caption={Generation-gated detail loading.}]
GenerationCounter generations = new GenerationCounter();

void loadDetails(CustomerId id) {
    GenerationCounter.Generation generation = generations.advance();
    tasks.submit(
            cancellation -> repository.loadDetails(id, cancellation),
            TaskCallbacks.onSuccess(details ->
                    generations.tryAccept(generation, details,
                            selectedDetails::setValue)));
}
\end{lstlisting}

If the user selects another customer before the first load finishes, the older
generation becomes stale. Its result is ignored.

\textbf{Failure policy.}

Failure should be modeled, not merely printed. Some failures belong in status
text. Some belong in a dialog. Some belong in a problem set. Some should be
logged and retried. The task callback should translate a technical failure into
the screen's presentation policy.

Avoid swallowing failures in background threads. A task service that hides
exceptions creates silent data loss. At minimum, the failure callback should
record the exception in diagnostics and set user-visible state.

\textbf{Cancellation policy.}

Cancellation should leave the screen in a coherent state. If a user cancels a
load, do old rows remain visible? Is the table cleared? Does status say
"Cancelled" or return to "Ready"? Is the command re-enabled? These are workflow
choices, not implementation accidents.

Cancellation is also not always user-visible. Closing a dialog may cancel a task
silently because the result no longer has an owner. In that case, stale-result
suppression and lifecycle cleanup are more important than a message.

\section{Worked screen: search, detail load, save, and export}

Consider a customer workspace with four asynchronous operations. A search field
loads table rows. Selecting a row loads detail. Editing the detail enables
Save. A toolbar command exports the current filtered table. This screen is
ordinary business Swing, but it contains most background-work problems. Search
results can arrive out of order. Detail loads can arrive after selection
changed. Save should block conflicting edits but not necessarily block search.
Export can run for seconds and should have cancellation and progress. A single
busy flag cannot describe this screen honestly.

The search operation is a request stream. Each new query supersedes older
queries. The presenter should keep a \api{GenerationCounter} for search, cancel
the previous search handle when possible, and accept only the newest result.
The table area may show a small busy mark, but the search field itself should
usually remain editable. The user should be able to type a better query while
the older query is still running. Disabling the search field for every request
would make the application feel slower than the service actually is.

The table rows should be replaced on the EDT in one batched operation. The task
body may call a repository, parse data, and create row objects. The callback
should check the generation and then update the observable list or table model
with a concise event. If the search result is stale, the callback should do
nothing except perhaps record a diagnostic count of ignored stale results. The
status line can say "Searching..." while the current handle is busy and "240
customers" after the accepted result. It should not display "0 customers" from
a stale empty result after a newer non-empty query has already won.

The detail load is another request stream, but its owner is the current
selection. Selecting customer A starts generation 1. Selecting customer B starts
generation 2. If A finishes after B, A must not replace the detail panel. If
the user clears selection, the detail generation should be invalidated so
queued callbacks cannot repopulate the panel. The detail busy value may block
only the detail form while leaving the table usable. Alternatively, the form may
show old detail until new detail arrives, with a small "loading selected
customer" mark. The product should decide; the task code should implement that
decision explicitly.

The save operation is not a request stream in the same way. A save usually
represents the current edit snapshot and should prevent conflicting edits until
the result is known. The detail form may be blocked by a \api{JLayer} overlay.
The Save command should be disabled while save is busy, while the form is clean,
or while validation has errors. The Cancel Edit command may or may not remain
available during save depending on transaction policy. The task body should
send a stable snapshot to the repository, not read live Swing fields while the
worker runs. The success callback should update dirty state and perhaps replace
the model with the saved version.

The export operation has a third shape. It may not change the table model at
all, but it may need progress and cancellation. Export busy should disable the
Export command and perhaps enable Cancel Export. It should not block typing in
the search field unless the export reads a live model that must stay frozen.
A safer design is to snapshot the export input on the EDT, then run export from
that snapshot in the task body. The user can continue working while export
uses the snapshot. The status text can say "Exporting 240 customers..." and
then "Export written" or "Export cancelled." The old table remains visible.

\begin{lstlisting}[caption={Four operations, four busy policies.}]
record WorkspaceBusy(
        ReadableValue<Boolean> searchBusy,
        ReadableValue<Boolean> detailBusy,
        ReadableValue<Boolean> saveBusy,
        ReadableValue<Boolean> exportBusy) {

    ReadableValue<Boolean> detailBlocked() {
        return saveBusy;          // Policy: save blocks detail editing.
    }

    ReadableValue<Boolean> tableBlocked() {
        return searchBusy;        // Policy: search blocks only table input.
    }
}
\end{lstlisting}

The record is illustrative rather than a required API. The point is that the
screen names separate busy facts. The table overlay observes search busy. The
detail overlay observes save busy or detail busy according to product policy.
The Save command observes dirty, valid, and save busy. The Export command
observes export busy and whether there are rows to export. A global application
status bar may observe a derived "any workspace task is busy" value while still
displaying the specific operation message.

The command derivations should be visible. A save command that is disabled
because the form is clean should have a different disabled reason from a save
command disabled because a save task is running. That difference can appear in
a composed tooltip or status message. It also matters for tests. A test that
sets dirty true, problems empty, and save busy false should see Save enabled.
A test that sets save busy true should see Save disabled with the busy reason.
Without structured state, these distinctions are often flattened into
\api{button.setEnabled(false)} with no explanation.

The overlay boundaries should match screen regions. The search overlay wraps
the table scroll pane, not the whole frame. The detail overlay wraps the form
body, not the toolbar. The export progress may live in the status area or a
small modeless progress panel. If every operation uses one root overlay, the
user loses useful parallelism. If no operation uses an overlay, incompatible
input may slip through. Region-level busy values are the practical middle
ground.

Progress messages should be operation-specific and stable. Search may not need
a progress bar because remote search has unknown duration. Detail load may need
only a short "Loading details" message. Save may need a blocking state and an
error path. Export may need determinate progress because the row count is known.
The screen should not force every task into the same visual component. A
Boolean busy overlay, a status label, and a structured progress model are
different tools.

Failures should preserve old state unless the workflow says otherwise. If
search fails, old rows may remain with a status message "Search failed; showing
previous results." If detail load fails, the detail panel may show a problem
state for the selected customer rather than old detail from another customer.
If save fails, dirty state remains and the user can retry. If export fails, the
table is unchanged and the export status reports the failure. Treating every
failure as "clear everything and show error" is rarely humane.

Cancellation policy differs by operation. Cancelling search simply means the
current query will not publish results. Cancelling detail load may leave the
detail panel empty or retain old detail depending on selection policy.
Cancelling save may not be offered after the request has reached a point of no
return. Cancelling export is often straightforward because export writes a file
or stream under local control. The UI should expose cancellation only where the
operation can honor it coherently. A Cancel button that cannot change the
outcome teaches users not to trust the screen.

The view lifecycle closes local operations. If the workspace tab closes, search
and detail handles should be cancelled or generation-invalidated. Save may need
a confirmation before closing if a write is in progress. Export might continue
as an application-level task if the product promises background export, or it
might be cancelled with the tab. The correct owner is a product decision, but
the implementation should make the owner visible: tab scope, dialog scope,
presenter scope, or application service.

Testing this workspace should not require a real database. Fake repositories
can expose latches or futures for search, detail, save, and export. A test can
start search A, start search B, complete B, then complete A and assert only B
is visible. Another test can start save, assert the detail overlay is busy and
Save disabled, then complete failure and assert dirty state remains. Another
can start export and assert search remains usable. These tests prove the busy
policy, not the repository.

The same scenario also explains why examples should use FlatLaf and MigLayout
without becoming decorative. The table, detail form, toolbar, and status area
are ordinary Swing surfaces. MigLayout gives dense, predictable arrangement.
FlatLaf gives modern component metrics and disabled-state rendering. Corusco
adds the state and lifecycle relationships. The screenshots should show
the combined result: not just a spinner, but a screen whose busy state is
visibly and behaviorally coherent.

The important lesson is that background work is not one feature. It is a set of
contracts around each operation. Search uses generation and partial blocking.
Detail load uses generation and selection ownership. Save uses a stable edit
snapshot and stronger blocking. Export uses progress and perhaps application
scope. The programmer's job is to name these contracts before the first worker
thread is started. Once they are named, Corusco provides small tools to express
them without turning the screen into listener bookkeeping.

\section{Plain Swing baseline}

Plain Swing has workable tools:

\begin{lstlisting}[caption={SwingWorker baseline.}]
SwingWorker<List<CustomerRow>, Void> worker =
        new SwingWorker<>() {
            @Override
            protected List<CustomerRow> doInBackground() {
                return repository.loadCustomers();
            }

            @Override
            protected void done() {
                if (!isCancelled()) {
                    List<CustomerRow> loaded = get();
                    tableRows.batch(list -> {
                        list.clear();
                        loaded.forEach(list::add);
                    });
                }
            }
        };
worker.execute();
\end{lstlisting}

The missing pieces are usually not the worker itself. They are task ownership,
busy state, cancellation policy, stale results, command enablement, overlay
binding, and cleanup. Corusco names those pieces.

This is why a plain Swing baseline remains worth teaching. \api{SwingWorker}
is not foolish; it captures the important idea that slow work should leave the
EDT. The problem is that many applications stop there. The worker becomes an
anonymous object inside a listener, and the rest of the screen has no stable
way to observe the operation. A task handle, a busy value, a cancellation token,
and a generation counter are not replacements for doing work. They are the
vocabulary around the work that lets the rest of the screen stay coherent.

\textbf{From plain Swing to Corusco.}

The Corusco shape for a busy screen is:

\begin{itemize}
\item Create a lifecycle-owned \api{TaskService}.
\item Submit tasks through the service.
\item Register task handles or the service with the view scope.
\item Bind busy state to commands, status, progress, and overlays.
\item Check cancellation in task bodies.
\item Use generation counters for request streams.
\item Keep callbacks short and EDT-safe.
\end{itemize}

This pattern works with virtual threads because virtual threads are an execution
choice, not a UI ownership model.

Migration from old code usually begins by extracting the slow call into a task
body that returns data. The next step is to move the successful UI update into
a callback that is known to run on the EDT. Then name the busy fact and bind it
to the button or command that started the operation. Only after that should the
screen grow a cancel button, progress display, busy overlay, or generation
guard. Trying to add every concern at once often recreates the original
listener tangle under different class names.

For a master-detail screen, generation counters should be added early. They are
small and prevent the worst class of asynchronous bug: an old result replacing
newer state without any exception. For a dialog, lifecycle cancellation should
be added early. A result that arrives after its dialog has closed no longer has
a legitimate owner. In both cases Corusco's task support makes the policy
explicit; it does not guess the policy for the application.

\begin{center}
\textbf{Busy review table.}
\end{center}

\noindent\begin{tabularx}{\linewidth}{@{}p{0.31\linewidth}X@{}}
\toprule
Observation & Review question\\
\midrule
Button can be clicked twice & Which busy or command state should block it?\\
Overlay shows but input still edits fields & Is overlay meant as input policy or decoration?\\
Old detail overwrites new selection & Where is stale-result suppression?\\
Cancel button does nothing visible & Does the task body check its cancellation token?\\
Status flickers rapidly & Should updates be coalesced or modeled as progress?\\
Failure disappears in logs only & Which user-visible failure state should be set?\\
\bottomrule
\end{tabularx}

\textbf{Worked example.}

\begin{examplebox}{Background Work}
\modernJavaExample{} demonstrates virtual-thread background work with explicit
EDT handoff. The Swing task examples in \api{corusco-examples} show busy overlay
binding and service-level busy state.
\end{examplebox}

\textbf{Review notes.}

When reviewing background work, trace the result path. Where is work started?
Where can it be cancelled? Where is busy state set? Which executor delivers
callbacks? Which line mutates Swing? What prevents an old result from winning?
If any answer is vague, the code is not done.

The review should also trace the user's perception. What becomes disabled? What
remains usable? What message appears during a long wait? What is shown after
success, cancellation, and failure? Does the screen make a second click
impossible, harmless, or meaningful? These questions are not cosmetic. They are
part of the concurrency design, because the user's next action is one of the
events racing with the background result.

A background chapter should also teach restraint. Not every computation belongs
in a task service. A small local calculation that always completes within a few
milliseconds can remain on the EDT, especially if moving it would complicate
state ownership. The boundary is unpredictability: I/O, remote calls, large
parsing, unbounded loops, and data sizes controlled by users. Put those outside
the EDT and make their results explicit.

Progress reporting should be designed at human scale. A worker may observe ten
thousand internal steps, but the user does not need ten thousand label updates.
Coalesce, throttle, or summarize progress before it reaches Swing. The model can
retain exact progress for diagnostics while the view displays a stable message
or a progress bar that changes at a readable cadence.

\section{Testing and operational review}

Background work is easy to test badly. A test that sleeps for a fixed time and
then looks at a label is neither fast nor reliable. A better test controls the
task boundary. Use latches or controlled executors to prove that the task body
runs off the EDT, the callback runs on the EDT, busy becomes false on the EDT,
and completion arrives after callback delivery. The existing Swing task tests
do exactly this: record whether the task body and callback are on the EDT, then
wait on the handle completion with a timeout. The timeout belongs to the test,
not to production UI code.

EDT delivery tests should be explicit because it is easy to accidentally create
a core task service with a non-Swing callback executor and then mutate Swing in
the callback. The test should submit a trivial task through
\api{SwingTaskServices.virtualThreads}, record
\api{SwingUtilities.isEventDispatchThread} in the task body and success
callback, and assert false then true. This is a small test with a large value:
it proves the factory used by presenters really supplies the Swing boundary.

Busy overlay tests should avoid manual screenshots as the first line of proof.
Construct a \api{JLayer}, install a \api{BusyOverlayLayerUI}, set busy true,
send a synthetic mouse or key event through \api{eventDispatched}, and assert
that the event is consumed. Paint into a small \api{BufferedImage} with an
opaque overlay color and assert a pixel. Install through
\api{BusyOverlayBinding}, close it, and assert the previous UI and cursor are
restored. These tests are mechanical, but they prove the visual/input policy
before a screenshot test ever runs.

Off-EDT busy updates deserve their own test. A busy value may be changed by a
worker. The binding should schedule the layer update onto the EDT and ignore
updates after close. A test can install the binding on the EDT, change the
value from a worker thread, drain the EDT, and assert the overlay state. Then
close on the EDT. This test protects the exact boundary that is most likely to
be broken by an enthusiastic optimization.

Cancellation tests should control the task body. Start a task that waits on a
latch, cancel the handle, release the latch, and assert that success does not
publish. For cooperative cancellation inside a loop, use a task body that
checks the token and records how far it ran. For repository-like code, mock or
fake a repository method that receives the token and verifies it was checked or
passed down. A Cancel button is only meaningful if cancellation reaches the
operation doing the work.

Stale-result tests should intentionally finish tasks out of order. Start a
first load and capture its generation. Start a second load and capture its
generation. Complete the second result first and assert it is accepted. Complete
the first result later and assert it is ignored. This test is more important
than a success-path test for master-detail screens because success order is the
bug. Without this test, a screen can pass every ordinary load test and still
overwrite current detail with old data in production.

Failure tests should assert both the callback path and the presentation state.
A task body throws. The failure callback runs on the EDT. Busy becomes false.
The command is re-enabled or remains disabled according to policy. The status
or problem state contains a useful message. The exception is available to
diagnostics. Do not assert only that an exception was thrown somewhere; the
user does not see a worker-thread stack trace. The user sees the screen state
after failure.

Close-path tests are the ones most likely to reveal leaks. Open a view, install
a busy overlay, submit a task, close the view, change the busy value, and drain
the EDT. The overlay should not reappear. The layer UI should be restored.
Model changes should not reach closed components. Outstanding handles should
be cancelled or generation-invalidated according to ownership. A test that
opens and closes the same screen several times should not accumulate listeners,
threads, overlays, or task callbacks.

Operational diagnostics should record task facts at the level that helps
support. A support dump for a busy screen can include current busy values,
active task names, generation number, last completion state, last failure class
and message, cancellation availability, and recent task durations. It should
not need to include thread dumps for ordinary "screen stayed busy" reports,
though thread dumps remain useful for deadlocks. The more task state is modeled,
the easier this dump becomes.

Logging should be restrained. Logging every progress update can become its own
performance problem. Log task start, terminal outcome, duration, cancellation,
and failure details. Sample or aggregate high-frequency progress. Include a
stable operation name or command id. If task code crosses a remote service or
database boundary, include enough correlation data to connect UI symptoms to
backend traces without leaking sensitive data into the UI log.

Reviewers should look for thread ownership by reading from left to right. The
command handler submits work. The task body uses no Swing components. The
callback updates models or components on the EDT. Busy values are derived from
service or handle state. Overlays and commands observe those values. Close
paths cancel or detach. If a component reference appears inside the task body,
pause. If a callback performs I/O, pause. If a worker thread mutates an
observable model already bound to Swing without an EDT boundary, pause.

The test pyramid for this chapter is therefore practical. Core task tests prove
service busy, cancellation, callback exclusivity, and generation counters
without Swing. Swing task tests prove EDT delivery and overlay binding.
Presenter tests prove command enablement, cancellation policy, and stale-result
suppression. Screenshot tests prove visible overlay and status composition.
Manual testing then checks feel: whether the chosen overlay boundary is too
broad, whether cancellation feels prompt, and whether progress text is useful.

This chapter also has a performance review. Count how many tasks can run at
once. Count how many UI updates each task can publish. Count how large each
published batch can be. Count whether each update can repaint a whole table,
whole form, or small region. The concurrency design is not complete until the
update cadence and repaint consequences are understood. Otherwise the program
merely moves the bottleneck from "work on EDT" to "too many callbacks on EDT."

The human review is shorter. When the user starts work, does the screen say
what is happening? Can the user cancel if cancellation is promised? Does the
screen prevent incompatible edits? Does it leave compatible work available? If
work fails, does the user know what happened and what state remains? If work is
cancelled, is the result coherent? If the answer to any of these is "it depends
on timing," the model is not explicit enough.

\section{Exercises}

\begin{exercisebox}
\begin{enumerate}
\item Convert a blocking refresh listener into a \api{TaskService} submission
with an EDT callback.
\item Bind service busy state to a \api{BusyOverlayBinding} and command
enablement.
\item Add a cancel button that cancels only the current task handle.
\item Add a generation counter to a master-detail load path.
\item Write a test that closes an overlay binding and verifies late busy updates
do not reinstall the overlay.
\item Decide cancellation policy for a load operation: keep old data, clear old
data, or show cancelled state.
\end{enumerate}
\end{exercisebox}

\textbf{Checklist.}

\begin{itemize}
\item Model busy state explicitly.
\item Disable or guard commands during incompatible work.
\item Deliver UI mutations on the EDT.
\item Keep task callbacks short.
\item Cancel or suppress stale completions.
\item Close task handles and services when their owner closes.
\item Use busy overlays as input policy, not merely decoration.
\item Test stale success, cancellation, and close paths.
\end{itemize}
